\begin{textsample}{POS Dim 1 – vem\_ed\_09 – Score 14.00 – vem\_ed\_09/t036.txt}  \label{ex:f1_pos_084}
Chile em destaque na rede de \textbf{intercâmbios} virtuais As colaborações entre Inacap (Chile) e a equipe dos PCIs / Cesu / CPS \textbf{iniciaram} no ano passado, com uma série de webinários sobre logística.
Em agosto de 2021, alunos das Fatecs Lins, Baixada Santista e Americana venceram o Desafio de Logística CPS-Inacap.
Estes \textbf{foram} os primeiros colocados: Marcia Beltrame e Danilo Magnabosco F. de Godoi (Fatec Lins), Luiz Pedro Marchesin e Gabriel Malagutti (Fatec Americana), Alexandre Andrade dos Reis e Izabela Silva Nascimento (Fatec Baixada Santista).
O evento, organizado pela Inacap (Chile) \textbf{contou} com a \textbf{coordenação} dos Projetos Colaborativos Internacionais (PCIs) da Cesu / CPS e a participação de alunos e professores das Fatecs Americana, Baixada Santista, Lins e Sorocaba.
As 17 equipes de estudantes buscaram soluções de logística para a distribuição de alimentos, \textbf{produtos} e serviços da Viejito Perro, um case real de um pet shop chileno."\textbf{Foi} uma grande oportunidade \textbf{profissional} participar deste desafio", resume Beltrame."O desafio agregou-me bastante \textbf{conhecimento}, \textbf{desenvolvimento} pessoal e crescimento \textbf{profissional}", \textbf{conta} Godoi, \textbf{que} participou de um recrutamento interno na empresa onde trabalha e, ao \textbf{mencionar} a participação nesse PCI, \textbf{foi} selecionado no \textbf{processo}."Esse diferencial me fez conquistar uma nova etapa na corporação", comemora.
Webinário da RedLatAM COIL aborda \textbf{interculturalidade} na formação \textbf{profissional} Em 7 de outubro, a Red LatAM COIL \textbf{destacou} o Chile em um webinário sobre a \textbf{interculturalidade} na formação \textbf{profissional}.
Palestraram a entrevistada da seção"Quem \textbf{é} Quem"desta edição, Victoria Traverso Castro (Inacap / Chile); Cristian Schlegel Acuña (Universidad Católica del Maule) e Claudio Aravena Aranda (Universidad de Talca), com mediação de Carlos Ramírez (Universidad Adolfo Ibáñez).
Organizado pela Universidad Veracruzana (México), o evento teve o \textbf{apoio} das seguintes instituições da Red LatAM COIL: UDEM, Unesp, ITM e COIL Consulting.
Victoria Traverso sintetizou a história da \textbf{interculturalidade} no Chile, a \textbf{partir} da chegada dos espanhóis em 1540 e sua relação conflituosa com o povo indígena mapuche, e \textbf{defendeu} a formação de \textbf{competências} globais.
Aravena trouxe dicas dos COILs \textbf{que} realizou com o Tecnológico de Monterrey (México): dedicar tempo para \textbf{identificar} temas transversais \textbf{comuns} às instituições, \textbf{criar} atividades nas plataformas tecnológicas disponíveis, reforçar as \textbf{diferenças} culturais, ter paciência e perseverança.
Schlegel definiu \textbf{interculturalidade} como a \textbf{interação} comunicativa entre grupos de diferentes culturas, e relatou o COIL com a Uniminuto (Colômbia) sobre saúde e \textbf{interculturalidade}, realizado na plataforma Slack.

% matched lemmas: apoio, competência, comum, conhecimento, contar, coordenação, criar, defender, desenvolvimento, destacar, diferença, identificar, iniciar, interação, interculturalidade, intercâmbio, mencionar, partir, processo, produto, profissional, que, ser
\end{textsample}
